
\documentclass{article}
\usepackage{colortbl}
\usepackage{makecell}
\usepackage{multirow}
\usepackage{supertabular}

\begin{document}

\newcounter{utterance}

\centering \large Interaction Transcript for game `guesswhat', experiment `Level\_3', episode 9 with reka{-}flash{-}3.1{-}t0.0.
\vspace{24pt}

{ \footnotesize  \setcounter{utterance}{1}
\setlength{\tabcolsep}{0pt}
\begin{supertabular}{c@{$\;$}|p{.15\linewidth}@{}p{.15\linewidth}p{.15\linewidth}p{.15\linewidth}p{.15\linewidth}p{.15\linewidth}}
    \# & $\;$A & \multicolumn{4}{c}{Game Master} & $\;\:$B\\
    \hline

    \theutterance \stepcounter{utterance}  
    & & \multicolumn{4}{p{0.6\linewidth}}{
        \cellcolor[rgb]{0.9,0.9,0.9}{
            \makecell[{{p{\linewidth}}}]{
                \texttt{\tiny{[P1$\langle$GM]}}
                \texttt{Let's play a guessing game. Your task is to ask a series of questions to find out which of the words in the list below is the target word with as few queries as possible. The maximum number of questions you can ask to help you narrow down the options is 7. Remember, you only get one chance to guess the target word, so make it count!} \\
\\ 
\texttt{The rules of the game are: } \\
\\ 
\texttt{1. Only ask one question per turn. DO NOT ASK MORE THAN ONE QUESTION.} \\
\\ 
\texttt{2. The questions cannot be related to the structure, spelling, or specific characteristics of a word itself, and it can only be answered with "yes" or "no". } \\
\\ 
\texttt{3. Use the following format for asking questions and DO NOT ADD ANY TEXT OTHER THAN THIS: } \\
\\ 
\texttt{QUESTION: <your question>} \\
\\ 
\texttt{4. You have only one opportunity to guess the target word. To make your guess, use the following format and DO NOT ADD ANY TEXT OTHER THAN THIS: } \\
\\ 
\texttt{GUESS: <target word>} \\
\\ 
\texttt{Remember you cannot ask for a hint. Make your question brief. Let’s start!} \\
\\ 
\texttt{List of candidate words: {[}'Halibut', 'Flounder', 'Carp', 'Catfish', 'Grouper', 'Mackerel', 'Minnow', 'Goldfish'{]}} \\
            }
        }
    }
    & \\ \\

    \theutterance \stepcounter{utterance}  
    & \multicolumn{4}{p{0.6\linewidth}}{
        \cellcolor[rgb]{0.9,0.9,0.9}{
            \makecell[{{p{\linewidth}}}]{
                \texttt{\tiny{[P1$\rangle$GM]}}
                \texttt{<reasoning>} \\
\texttt{Let's see. I need to figure out which fish is the target word from the list provided. The list is {[}'Halibut', 'Flounder', 'Carp', 'Catfish', 'Grouper', 'Mackerel', 'Minnow', 'Goldfish'{]}. The user wants me to ask yes/no questions to narrow it down in 7 steps, then guess.} \\
\\ 
\texttt{First, I should think about common characteristics of these fish. Maybe start with something broad. Let's see. Are all these fish saltwater or freshwater? For example, Halibut and Flounder are typically saltwater, while Carp and Catfish are freshwater. Grouper, Mackerel, Minnow, and Goldfish can vary, but Goldfish are definitely freshwater. So maybe the first question is about saltwater vs freshwater. That could split the list into two groups. Let's check: if the answer is yes (saltwater), then possible candidates are Halibut, Flounder, Grouper, Mackerel, Minnow. Wait, Minnows can be both, but maybe the common ones are freshwater. Wait, Minnows are usually freshwater. So if the answer is yes (saltwater), then Halibut, Flounder, Grouper, Mackerel. If no, then Carp, Catfish, Minnow, Goldfish. That's a good split. Let's go with that.} \\
\\ 
\texttt{So first question: "Is the target word a saltwater fish?" Let's see. If yes, then the remaining options are Halibut, Flounder, Grouper, Mackerel. If no, then Carp, Catfish, Minnow, Goldfish. Let's assume the answer is yes. Then next step. Maybe ask about size? Halibut is a large fish, Flounder is flat, Grouper is large, Mackerel is oily. Alternatively, maybe ask about common aquarium fish. Goldfish are common, but others like Carp are sometimes kept, but maybe not as much. Wait, the list includes Goldfish, which is definitely an aquarium fish. So if the answer to the first question is no (freshwater), then Goldfish is a candidate. But if the answer is yes, then maybe not. Wait, but the first question is about saltwater. So if the target is Goldfish, it's freshwater, so the first answer would be no. So} \\
            }
        }
    }
    & & \\ \\

    \theutterance \stepcounter{utterance}  
    & & & \multicolumn{2}{p{0.3\linewidth}}{
        \cellcolor[rgb]{0.9,0.9,0.9}{
            \makecell[{{p{\linewidth}}}]{
                \texttt{\tiny{[GM$|$GM]}}
                \texttt{Invalid format. Guesser must use the form 'QUESTION: ' or 'GUESS: '.} \\
            }
        }
    }
    & & \\ \\

    \theutterance \stepcounter{utterance}  
    & & & \multicolumn{2}{p{0.3\linewidth}}{
        \cellcolor[rgb]{0.9,0.9,0.9}{
            \makecell[{{p{\linewidth}}}]{
                \texttt{\tiny{[GM$|$GM]}}
                \texttt{abort game} \\
            }
        }
    }
    & & \\ \\

\end{supertabular}
}

\end{document}
